% STATUS: ACTIVE -- companion reference notes (appendix) for formB_ws.tex.
%          The main file is equation-only; every convention, rationale,
%          warning and provenance lives here, under the SAME section numbers.
% ======================================================
% Compile: xelatex formB_ws_ref.tex
% Path:    reference/eq17_analysis/derivation/formB_ws_ref.tex
% Main:    formB_ws.tex
% ======================================================

\documentclass[11pt,a4paper,fontset=none]{ctexart}

\usepackage[fleqn]{amsmath}
\usepackage{amssymb,bm}
\setlength{\mathindent}{0pt}
\usepackage{geometry}
\geometry{margin=2.0cm}
\usepackage{array,booktabs}

\setlength{\parindent}{0pt}
\setlength{\parskip}{0.80em}
\linespread{1.08}
\setlength{\emergencystretch}{2em}

\begin{document}

\section*{Conventions (apply to the whole of formB\_ws.tex)}

\textbf{Normalization --- the derivation is fully dimensionless; the bar
marks the normalized version of a quantity that has a physical twin.}
Lengths by the particle radius $R$: $\bar{w} = w/R$ is the probe-center
height on the wall-normal axis (the same axis the companion documents call
$\bar{h}$), $\bar{w}_s$ is the contact (surface) position, nominal $1$; the
unknown is the offset $\bar{w}_s - 1$.  Gain by
$a_o = \Delta t/(\gamma_N R)$: $\bar{a} = a/a_o \in [0,1]$ (exactly the
quantity of the companion note, saturating at $1$); the unbarred
$a_w = a_o\,\bar{a}_w$ is the physical gain, $1/\mathrm{pN}$.  Forces by
$a_o$: $\bar{f} = a_o f$ is the normalized far-field one-step displacement
of the force $f$, so $\bar{a}\,\bar{f} = a\,f$ identically and every
equation of the derivation is dimensionless.  The parameters $b$ (a length
in units of $R$, nominal $9/8$) and $p$, and the noise symbols
$\varepsilon_w$, $n_w$, $n_a$, follow the companion documents' unbarred
writing; their normalization is declared here once ($n_w$: position noise
in units of $R$; $n_a$: readout noise on the \emph{normalized} gain,
i.e.\ the house doc's physical readout noise divided by $a_o$).
Background symbols: $\gamma_N = 6\pi\eta R$ (far-field Stokes drag);
$\gamma(\bar{w}) = \gamma_N\,c(\bar{w})$ with $c(\bar{w})$ the wall-effect
correction factor, unknown at run time --- it appears only inside the
$\bar{a}_w$ identity in S8, never in a run-time quantity;
$\delta\bar{w}_r$ = closed-loop regulation residual (the house doc's
$\delta\bar{h}_r$ renamed).  Estimates wear the hat over the
barred symbol: $\hat{\bar{a}}_w$, $\hat{b}_w$, $\hat{p}_w$,
$\hat{\bar{w}}_s$.  $a_o$ appears only at the two interfaces: force conversion
$\bar{f} = a_o f$ and the gain readout $\bar{a}_{wm} = a_{wm}/a_o$.
Discrete index $k$, sample time $\Delta t$.
(Correction log: an earlier draft wrote ``$\bar{a}_w$ in $1/\mathrm{pN}$''
--- wrong, fixed 2026-07-31; a second draft defined $e_{a_w}$ on the
physical gain --- superseded 2026-07-31 by the fully normalized scheme
below.  Both swaps are value-preserving: $\bar{a}\bar{f} = af$.)

\textbf{Error convention.}  $e_{\theta} = \theta - \hat{\theta}$ (true minus
estimate) for every estimated quantity.  The gain error is
$e_{\bar{a}_w} = \bar{a}_w - \hat{\bar{a}}_w$, dimensionless; with $F_{dw}$
built from normalized efforts $\bar{f}_{dw}$, the product
$F_{dw}\,e_{\bar{a}_w}$ is a normalized length, consistent with the
$\delta\bar{w}$ recursion.  This matches the companion note exactly: its
$\bar{a}$ saturates at $1$ and multiplies $F_{dw}$ directly, which is
dimensionally consistent only under the normalized-force reading.

\textbf{Notation rule} (user, 2026-07-31): only symbols already defined in
the two companion documents are used --- ``Estimation of motion gain\_R''
and \texttt{5state\_expgain\_hd}; every addition needs prior discussion.
Expressions are kept fully expanded --- no substitution variables (house
rule of \texttt{aprime\_jacobian\_AB.tex}).

\textbf{Cross-document symbol table.}
\begin{center}
\begin{tabular}{lll}
\toprule
symbol & meaning in formB\_ws & clash to avoid \\
\midrule
$b$ & Form B length scale, units of $R$, nominal $9/8$ &
  exponent $b$ of \texttt{5state\_expgain\_hd} \\
$p$ & Form B exponent, nominal $1$ &
  power $p$ of \texttt{5state\_powerlaw\_hd} \\
$\bar{w}$, $\bar{w}_s$ & height / contact position &
  $\bar{w} \equiv \bar{h}$ of the expgain and powerlaw files \\
$\bar{a}$ & normalized gain $a/a_o$ &
  unbarred $a_w$ is the physical gain, $1/\mathrm{pN}$ \\
\bottomrule
\end{tabular}
\end{center}

\textbf{Imported sources.}  Base chain (control law $\to$ boxed motion error)
from \texttt{4state\_del\_hd.tex} / \texttt{5state\_expgain\_hd.tex},
imported in S3 with the $\bar{h}\to\bar{w}$ rename (map in S3 notes below);
shape selection and Jacobian verification from
\texttt{gainlaw\_shape\_selection.tex} and \texttt{aprime\_jacobian\_AB.tex};
truth curve from \texttt{two\_sphere\_truth\_source.md}.

\section*{S1 notes --- Gain law}

\textbf{Representation.}  The writing
$\bar{a} = 1 - (1 + (\bar{w}-\bar{w}_s)/b)^{-p}$ was fixed by the user
(2026-07-31).  It is algebraically identical to the length writing
$1 - [\,b/((\bar{w}-\bar{w}_s)+b)\,]^{p}$ of
\texttt{gainlaw\_shape\_selection.tex}; the equivalence, the three parameter
Jacobians in this writing, and the wall-shift identity
$\partial\bar{a}/\partial\bar{w}_s = -\bar{a}'$ were all checked symbolically
(MATLAB, residual exactly $0$, 2026-07-31), so every verified asset of
\texttt{aprime\_jacobian\_AB.tex} carries over by rewriting alone.

\textbf{Contact property and the amplitude-writing trap.}  At
$\bar{w}=\bar{w}_s$ the base is $1$, so $\bar{a}=0$ for any $(b,p)$:
$\bar{w}_s$ is the contact position by construction, and the unknown offset
$\bar{w}_s-1$ is what the estimator will carry.  The bare amplitude writing
$1-\lambda_{amp}(\bar{w}-\bar{w}_s)^{-p}$ loses this: fitted against the
truth it returns $\bar{w}_s=-0.063$, which is not a surface position, and it
drags fractional powers and $\ln$ terms into every Jacobian.  Do not use it.
(Audit note in \texttt{gainlaw\_shape\_selection.tex}.)

\textbf{Limits and anchors} (removed from the main file 2026-07-31, kept
here).  Near wall $\bar{a} = (p/b)(\bar{w}-\bar{w}_s) +
O((\bar{w}-\bar{w}_s)^2)$; far field $1-\bar{a} \to
(b/(\bar{w}-\bar{w}_s))^{p}$.  Both anchors are derived, not fitted
(truth = two-sphere $\perp$, Jeffrey \& Onishi 1984): the far-field truth
$1-\bar{a}\sim(9/8)/(\bar{w}-\bar{w}_s)$ forces $p=1$ and $b=9/8$ (the
reflection coefficient); the near-wall truth slope is $1$ while the model
slope at the anchored pair is $p/b = 8/9$ --- the truth sits $9/8 = 1.125$
above the model ($+12.5\%$ of the model; equivalently the model is
$11.1\%$ below the truth).  The two ends cannot be satisfied
simultaneously: this tension (quoted as $12.5\%$, direction as stated) is
the quantified shape imperfection, and it funds the parameter prior
widths in S10 (the same role the interpolation-gap $\sup$ plays for the
expgain and powerlaw priors).

\textbf{Fit record} (minimax over $\bar{w}\in[1.1,10]$, $2\times10^4$
samples): three parameters free give $\sup$ error $0.429\%$ at
\begin{equation*}
(b,\ p,\ \bar{w}_s) = (1.0566,\ 0.9563,\ 0.9936);
\end{equation*}
$(p{=}1,\ b{=}9/8)$ pinned with $\bar{w}_s$ the only estimate: $0.779\%$.
Form A under the same protocol
(\texttt{gainlaw\_shape\_selection.tex}): $4.256\%$, about $10\times$
worse; its two asymptotic anchors have no joint solution.

\textbf{Why $a_o$ is fixed, not a state.}  Two independent reasons.
(i) \emph{Honesty}: the $y_2$ variance-readout chain currently carries the
open motion-time defect ($+4\sim5\%$, defect 2); a freed $\hat{a}_o$ acts as
a bias sponge for it --- measured on the expgain 7a prototype, $\hat{a}_o$
parks at $+1.2\%$ off truth while the overall RMS improves
($0.477\%\to0.144\%$).  That RMS comparison is exactly the no-trade-off
fingerprint: both numbers share the polluted readout chain, so it must not
decide the question.  Fixed until defect 2 is re-derived
(\texttt{Cdpmr\_Cn\_derivation.tex}).
(ii) \emph{No absorber exists}: the far-field limit of $\bar{a}$ is pinned at
$1$ for every $(b,p,\bar{w}_s)$, so no shape parameter can mimic an amplitude
error; an $a_o$ error is an irreducible multiplicative bias, uniform in
height ($\delta a_w/a_w = \delta a_o/a_o$ at every $\bar{w}$), bounded by
the calibration tolerance chain
$\delta a_o/a_o=\sqrt{(2\,\delta R/R)^2+(\delta\eta/\eta)^2}\approx 2.8\%$
($a_o = \Delta t/(\gamma_N R)$, $\gamma_N = 6\pi\eta R$).

\textbf{Acceptance consequences} (S11): inject $a_o$ perturbations of
$\pm2.8\%$ and measure the leakage into $\hat{b}_w$, $\hat{p}_w$,
$\hat{\bar{w}}_s$ (parameter-invariance style test); revisit the
fixed-vs-free question only after defect 2 closes.

\textbf{$p=1$ special case.}  With the far-side anchors in force the law and
all its Jacobians are rational (no \texttt{pow}/\texttt{exp}/\texttt{log}),
which the $10\,\mathrm{kHz}$ FPGA loop can evaluate directly.

\section*{S2 notes --- Gain rate and parameter Jacobians}

\textbf{Provenance and verification.}  All three Jacobians are the Form B
set of \texttt{aprime\_jacobian\_AB.tex}, rewritten into the photo
representation; each scalar bracket equals
$\partial \ln \bar{a}'/\partial\theta$ (log-derivative factorization
$J_\theta = [\text{scalar}]\cdot\bar{a}'$), and each was re-verified
symbolically in this writing (MATLAB symbolic check, residual exactly $0$,
2026-07-31).  The acceptance protocol for any future
re-expansion is Richardson ratio $\approx 100$ with a mandatory negative
control collapsing to $\approx 10$, per coordinate direction
(\texttt{aprime\_jacobian\_AB.tex}).

\textbf{Menq-note comparison.}  The companion note's $J_b$, $J_p$ were
written for Form A and carry sign errors (the $-$ from
$d\,e^{-u}/d\theta$ was dropped; its $J_{w_s}$ is correct, which proves the
three were mutually inconsistent).  The form swap replaces that whole block;
nothing needs patching line by line.

\textbf{Expansion validity.}  First order in
$(e_{b_w}, e_{p_w}, e_{\bar{w}_s})$; the $O(\|e\|^2)$ remainder is dropped ---
declared here once so the main file can stay clean.  Both $\bar{a}'$
evaluations share the anchor point $\bar{w}_d-\hat{\bar{w}}_s$, so this
expansion covers parameter errors only; the evaluation-point error
(command vs true position, $O(\bar{a}''\delta\bar{w})$) is handled where the
process model is built (S4), not here.

\textbf{Identifiability warning} (for S10, not a reason to change the form):
all three $J$'s multiply the same scalar excitation in the error dynamics,
so per step the parameter subspace is excited along a single direction;
measured on the 3-parameter fit: Fisher condition $\sim$270,
$\cos(b,p$-directions$) = 0.979$, $\cos(b,w_s) = 0.642$.  Handled by priors
and trajectory design; individual $\hat{b}_w$, $\hat{p}_w$ must not be
quoted separately (four-rules note of the 2026-07-29 memory).

\section*{S3 notes --- True dynamics}

\textbf{Import and rename map.}  The boxed $\delta\bar{w}$ recursion and the
$F_{dw}$, $\varepsilon_w$ definitions are the boxed motion-error result of
the base chain (\texttt{4state\_del\_hd.tex} /
\texttt{5state\_expgain\_hd.tex} p.1--2), imported with the axis rename:
\begin{center}
\begin{tabular}{llllllll}
\toprule
there & $\bar{h}$ & $\delta\bar{h}$ & $\Delta\bar{h}_d$ & $F_{dh}$ &
  $\varepsilon_h$ & $n_h$ & $f_{dh}$ \\
here  & $\bar{w}$ & $\delta\bar{w}$ & $\Delta\bar{w}_d$ & $F_{dw}$ &
  $\varepsilon_w$ & $n_w$ & $\bar{f}_{dw}$ \\
\bottomrule
\end{tabular}
\end{center}
The rename is mechanical; the equations were not re-derived here, and their
SSOT remains the base-chain file.

\textbf{Imported premises} (carried by the boxed recursion; re-audit list):
(1) measurement delay exactly $d$ steps, no jitter;
(2) unmodeled disturbance $\delta\bar{w}_D \equiv 0$ (dropped from the
control law and hereafter);
(3) closed-loop pole $\lambda_c$ known exactly to the estimator;
(4) gain-error back-step: $a_w[k-i]-\hat{a}_w[k-i] \approx e_{a_w}[k]$ over
the $d$-step window --- the dropped term
$(\bar{w}_d[k]-\bar{w}_d[k-i])(a_w'-\hat{a}_w')$ is form-dependent and must
be re-bounded with Form B's $\bar{a}'$, $\bar{a}''$ before S11 signs off
(open TODO, tracked here).

\textbf{Delay line.}  The states $\delta\bar{w}_1,\dots,\delta\bar{w}_3$
follow the State-section convention of \texttt{5state\_expgain\_hd} (the
measurement carries $\delta\bar{w}[k-d]$, i.e.\ $\delta\bar{w}_1$); the
companion note has no delay-line states, which is why its error system is
not closed --- bringing the delay line into the true dynamics here is the
closure fix.  $d=2$ is the project's canonical setup; the sums stay generic
in $d$.

\textbf{Reading of the true-gain line.}  Written at $[k{+}1]$ in the
companion note's composition form
$\bar{a}_w[k+1] = \bar{a}_w(\bar{w}[k+1])$: the true gain is evaluated at
the true position $\bar{w} = \bar{w}_d-\delta\bar{w}_3$, i.e.\ it trails the
command by the tracking error, and this $[k{+}1]$ form is what the process
model (S4) Taylor-expands.  The estimator (S5) will evaluate at $\bar{w}_d$
with the estimated parameters, and the difference splits into a parameter
part (S2 Jacobians) and an evaluation-point part
($O(\bar{a}''\delta\bar{w}_3)$, bounded in S4).  $e_{a_w}$ appears inside
the true $\delta\bar{w}_3$ recursion (self-modulation: the true trajectory
depends on the estimation error through the control); this coupling is
inherited from the base chain, not introduced here.

\textbf{Constants.}  $\bar{w}_s$, $b_w$, $p_w$ carry $[k+1]=[k]$ with zero
process noise; the honesty condition that legitimizes $Q_\theta = 0$
($\sup|\theta_{\mathrm{eff}}-\theta_0| \lesssim \sqrt{P[0]}$) is stated and
tested in S10/S11, not assumed silently.

\section*{S4 notes --- Process model}

\textbf{Decision record} (user, 2026-07-31).  S4 follows the companion
note's structure: $\bar{a}_w$ is a state propagated by $\bar{a}'$, and
$b_w$, $p_w$ are \emph{not} pinned --- the full parameter set rides along,
so the closed system is 7 states:
$[\delta\bar{w}_1, \delta\bar{w}_2, \delta\bar{w}_3, \bar{a}_w, b_w,
p_w, \bar{w}_s]$ (parameter order $b_w, p_w, \bar{w}_s$ throughout, user
2026-07-31).  Identifiability of $(b_w, p_w, \bar{w}_s)$ is handled by
priors and trajectory design (S10); individual $\hat{b}_w$, $\hat{p}_w$
must not be quoted separately (S2 warning).

\textbf{First-order Taylor and its two dropped terms.}  The boxed row is
$\bar{a}_w[k+1] = \bar{a}_w(\bar{w}[k+1])$ expanded to first order about
$\bar{w}[k]$, with $\bar{a}'$ evaluated at the desired gap (the note's
highlighted line).  Dropped without statement in the note, declared here:
(i) the Taylor remainder $O(\bar{a}''\cdot(\text{increment})^2)$;
(ii) the evaluation-point error
$O(\bar{a}''\cdot\delta\bar{w}_3\cdot\text{increment})$ from using
$\bar{w}_d-\bar{w}_s$ instead of $\bar{w}[k]-\bar{w}_s$.  Closed form
\begin{equation*}
\bar{a}'' = -\frac{p(p+1)}{b^2}
  \Bigl(1+\frac{\bar{w}-\bar{w}_s}{b}\Bigr)^{-p-2},
\qquad
\sup_{\bar{w}>\bar{w}_s}|\bar{a}''| = \frac{p(p+1)}{b^2},
\end{equation*}
the sup attained at contact (nominal $(p,b)=(1,\tfrac98)$:
$\tfrac{128}{81}\approx 1.58$), so both bounds are computable once the
trajectory class is fixed (to be instantiated in S11).

\textbf{Full-state assembly.}  The section closes with the state vector
$\mathbf{x} = [\delta\bar{w}_1, \delta\bar{w}_2, \delta\bar{w}_3,
\bar{a}_w, b_w, p_w, \bar{w}_s]^{\top}$ ($d=2$) and the complete boxed
true-state block --- the house form of \texttt{5state\_expgain\_hd}'s
``True state equation''.  Rows 1--3 repeat the S3 box (deliberate, same as
the house doc); row 4 is the gain process model; rows 5--7 the constants.
The standalone gain row above the block is the derivation step (Taylor of
the S3 composition), left unboxed.

\textbf{Structural risk record (defect 1).}  Differential propagation
integrates a static curve as a time recursion.  On the expgain law this
produced the measured defect-1 signature: descent peak $11.03\%$
(differential) vs $5.05\%$ (algebraic evaluation); adding the $y_2$
measurement \emph{worsens} $\times 2.05$ vs \emph{improves} $\times 0.70$;
constants' observability $159$--$7886\times$ worse.  The mechanism is
structural and expected to transfer, but it has \emph{not} been measured on
Form B --- hence the S11 fingerprint criterion is mandatory, and the
algebraic alternative (evaluate $\hat{\bar{a}}_w$ from
$(\hat{\theta}, \bar{w}_d)$ each step; identical Jacobian machinery, one
state row differs) is kept on record here as the pivot plan.

\textbf{Observability note.}  With $\bar{a}_w$ a state, the constants
influence the filter only through the $\bar{a}$-row coupling
$J_\theta[k]\cdot(\text{increment})\cdot e_\theta$: during hold the
increment $\approx 0$ and $(b_w, p_w, \bar{w}_s)$ are unobservable (hold
desert); motion is required.  The exact channel
$\partial\bar{a}/\partial\bar{w}_s = -\bar{a}'$ makes $e_{\bar{w}_s}$ and
$e_{\bar{a}_w}$ instantaneously collinear in any gain readout; separation
comes only from height variation.  This drives the trajectory design of the
S11 injection test.

\textbf{Noise container preview (S7).}  The noise entering the
$\bar{a}$-row is $\bar{a}'\cdot\varepsilon_w$ $\to$ rank-one $Q$ block
across $(\delta\bar{w}_3, \bar{a}_w)$ (expgain pattern).  The Taylor and
evaluation-point truncations are coherent, trajectory-driven errors and
must \emph{not} be put into $Q$ (container rule); they are exactly what the
S11 fingerprint test watches.

\section*{S5 notes --- Estimator}

\textbf{Mirror construction.}  The estimator is the process model with every
quantity it cannot know removed and every remaining quantity replaced by its
estimate: $F_{dw}\,e_{\bar{a}_w}$ is dropped because $e_{\bar{a}_w}$ is
unknown (the estimator does know $F_{dw}$ --- it computes the control ---
but not the error it multiplies); $\varepsilon_w$ is dropped as zero-mean
noise.  $\delta\hat{\bar{w}}_3$ therefore contracts with the ideal
closed-loop pole $\lambda_c$ alone.  The gain row mirrors S4 with
$\hat{\bar{a}}'[k]$ evaluated at the estimates (the companion note's second
highlighted line): run-time computable from $\bar{w}_d$ and the state
estimates only --- $c$-free by construction.

\textbf{Prediction-only, by design.}  Following the companion note, S5
carries no correction terms; the gains enter with the measurement (S7) and
the filter equations (S9).  (\texttt{5state\_expgain\_hd} folds the
$\ell_{ij}$ update terms into its Estimator section instead --- a layout
difference only, noted to prevent a false mismatch reading.)

\textbf{Closure fix.}  The companion note has no
$\delta\hat{\bar{w}}_j$ states, so its $e_w = \delta\bar{w}-\delta\hat{\bar{w}}$
had no dynamics and the error system did not close.  With the hat delay
line present, $e_{\delta\bar{w}_j} = \delta\bar{w}_j - \delta\hat{\bar{w}}_j$
acquires the shift dynamics, and S6 (True $-$ Estimator) closes the
7-dimensional error system.

\section*{S6 notes --- Error dynamics}

\textbf{Derivation split (the section's single first-order declaration).}
The exact difference of the two gain rows decomposes as
\begin{equation*}
(\bar{a}' - \hat{\bar{a}}')\cdot
\bigl[\Delta\bar{w}_d + (1-\lambda_c)\delta\hat{\bar{w}}_3\bigr]
\;+\;
\bar{a}'\cdot\bigl[(1-\lambda_c)e_{\delta\bar{w}_3}
 + F_{dw}e_{\bar{a}_w} + \varepsilon_w\bigr].
\end{equation*}
The first factor is replaced by the S2 Jacobian expansion (evaluated at the
estimates); in the second group $\bar{a}' \to \hat{\bar{a}}'$, which
introduces only second-order terms
$O\bigl((\bar{a}'-\hat{\bar{a}}')\cdot(e,\varepsilon_w)\bigr)$
(error-times-error and error-times-noise, the latter landing in $q_4$'s
coefficient) --- dropped, declared here.  The row-wise result is
structurally identical to the companion note with its $e_w$ read as our
$e_{\delta\bar{w}_3}$.

\textbf{Prediction-only form.}  The boxed error state equation carries no
correction: $-\mathbf{L}[k]\mathbf{H}[k]\mathbf{e}[k]$ joins in S9 once the
measurement (S7) exists.  (\texttt{5state\_expgain\_hd} folds the
$\ell_{ij}$ terms into its row-wise error dynamics instead --- layout
difference only.)

\textbf{$\mathbf{q}$ structure and a known cross-correlation.}
$q_3 = -\varepsilon_w$, $q_4 = \hat{\bar{a}}'\varepsilon_w =
-\hat{\bar{a}}'q_3$: the gain block is rank one, the expgain pattern.
$\varepsilon_w$ contains $(1-\lambda_c)n_w$, so $\mathbf{q}$ correlates
with the $y_1$ measurement noise ($E[\mathbf{q}\,\mathbf{r}^{\top}] \neq 0$)
--- the same declared-open item as \texttt{5state\_expgain\_hd}'s NOTES;
to be handled or bounded in S8.

\textbf{Single-direction excitation, now visible.}  Columns 5--7 of
$\mathbf{F}_e$'s row 4 share the common factor
$[\Delta\bar{w}_d + (1-\lambda_c)\delta\hat{\bar{w}}_3]$: per step the
three parameter errors are excited along the single direction spanned by
$(J_{b_w}, J_{p_w}, J_{w_s})$ --- the S2 identifiability warning made
concrete.  Hold $\Rightarrow$ the factor $\approx 0$ $\Rightarrow$
parameter hold-desert (S4 note); this drives the S10 priors and the S11
trajectory design.

\section*{S7 notes --- Measurement and innovation}

\textbf{Readout provenance.}  $y_2 = \bar{a}_{wm} = a_{wm}/a_o$ is the
variance-readout gain: the physical $a_{wm}$ is reconstructed from
$\hat{\sigma}^2_{\delta\bar{w}_r}$ through the $C_{dpmr}/C_n$ chain, which
arrives in S8 together with its OPEN flag (defect 2: the chain assumes a
stationary, noise-only-driven loop; motion-time excess $+4\sim5\%$).  The
normalization by $a_o$ is the second and last place $a_o$ appears
(Conventions).

\textbf{Delay back-off --- house form by decision} (user, 2026-07-31).
The back-off uses the command displacement only,
$\bar{a}_w[k-d] = \bar{a}_w[k] - \bar{a}'\nabla_d\bar{w}_d$, matching
\texttt{5state\_expgain\_hd} exactly.  Declared truncations:
(i) the tracking-error difference over the delay window is dropped ---
the exact increment is $\nabla_d\bar{w}_d - \delta\bar{w}_3 +
\delta\bar{w}_1$, so the dropped term is
$\hat{\bar{a}}'(\delta\bar{w}_3 - \delta\bar{w}_1)$, zero in hold and
nonzero in motion transients (the house NOTES record this channel
$1.3\times$ low); (ii) the back-off Taylor is first order, remainder
$O(\bar{a}''\,\nabla_d\bar{w}_d^{\,2})$; (iii) $\bar{a}'$ in the back-off
is evaluated at the command point --- the same evaluation-point class
declared in S4, $O(\bar{a}''\,\delta\bar{w}_3\,\nabla_d\bar{w}_d)$.
Consequence for S8: since
$\mathbf{H}$ no longer carries the delay-window channel, $R_2$ must ---
the house doc's $+\,d\,Q_{44}$ term.  The exact-$\mathbf{H}$ alternative
($\mathbf{H}$ row 2 columns 1/3 $= \mp\hat{\bar{a}}'$) was derived and
symbolically verified (constrained $\delta\bar{w}_3$-dependence cancels
exactly; $\delta\bar{w}_1$-sensitivity matches the exact readout;
2026-07-31) and is kept on record here as an available upgrade --- it is
the 7b-analogue of 7a's $H(2,1) = -a'$ fix.

\textbf{Form B removes a house-doc truncation.}  In
\texttt{5state\_expgain\_hd}, $\partial y_2/\partial a_h = 1 +
\nabla_d\bar{h}_d\hat{B} \to 1$ was itself a truncation ($\sim0.4\%$).
Here $\bar{a}'$ is built from $(b_w, p_w, \bar{w}_s)$ and does not depend
on the state $\bar{a}_w$, so $\partial y_2/\partial\bar{a}_w = 1$ exactly
--- nothing is dropped.

\textbf{Whitening pointer.}  $y_2$ as written is the raw EWMA readout; it
is AR(1), and treating it as white overestimates its information
$39\times$ (measured; $a_{cov}$-invariance FAIL $3.14\%$ raw vs PASS
whitened).  The whitening transform is introduced together with $R_2$ in
S8, modifying $(y_2, \mathbf{H}\ \text{row 2}, R_2)$ as one unit ---
S8 discussion item.

\textbf{First-order declaration.}  In $e_{y_2}$ the mixed second-order
terms $O\bigl((\bar{a}'-\hat{\bar{a}}')\cdot e_{\delta\bar{w}}\bigr)$ are
dropped, same one-declaration convention as S6.

\section*{S8 notes --- Process and measurement noise}

\textbf{Declared exclusions (house-doc lineage).}
(i) The past thermal steps
$(1-\lambda_c)\sum \bar{a}_w[k-i]\bar{f}_T[k-i]$ are cross-step correlated
and stay out of $Q$ --- the container is knowingly incomplete, the same
declaration the house doc makes.
(ii) $\varepsilon_w$ contains $(1-\lambda_c)n_w$, so
$E[q_3\,r_1] = -(1-\lambda_c)\sigma^2_{n_w} \neq 0$ is dropped (the
standard-KF independence assumption is kept).
(iii) Run time evaluates every $Q$ entry at the current estimate ---
stated once here, fixing the evaluation-point asymmetry the house doc
left undeclared.
(iv) $d\,Q_{44}$ in $R_2$ carries the delay-window gain wander because
S7's house back-off leaves that channel out of $\mathbf{H}$ (S7 notes;
the house NOTES record this channel $1.3\times$ low).

\textbf{OPEN --- defect 2 lives in the $\sigma^2_{\delta\bar{w}_r}$
formula.}  $C_{dpmr}$, $C_n$ are derived
(\texttt{Cdpmr\_Cn\_derivation.tex}) for a \emph{stationary,
noise-only-driven} closed loop.  Measured (2026-07-28): during commanded
motion $\mathrm{Var}(\delta\bar{w}_r)$ exceeds the formula by $+4\sim5\%$,
growing with command frequency ($0.25/0.5/1$\,Hz $\to$
$+0.9/+2.4/+5.2\%$) and with the $a_{pd}$ EWMA time constant; stationary
excess is only $+1.0\%$; seven falsifications excluded every other
suspect.  The inversion therefore biases $y_2$ high by the same amount
during motion, and everything downstream of $y_2$ inherits it (the
$+1.2\%$ parked $\hat{a}_o$; the polluted B-layer numbers; the blocked
freeze gate).  Re-deriving \texttt{Cdpmr\_Cn\_derivation.tex} with motion
included is the mainline's first priority.

\textbf{The correct (whitened) $y_2$ treatment --- production default.}
$\hat{\sigma}^2$ is an EWMA, so $y_2$ is AR(1) with pole $(1-a_{cov})$:
$K_{var}$ fixes the per-sample \emph{variance} but not the
sample-to-sample \emph{correlation}.  Fed raw at every step, the filter
overcounts information by $(2-a_{cov})/a_{cov} \approx 39\times$ (exact
analytic factor; the 2026-07-27 incident --- the fake $0.92$--$0.95$
results).  Whitening feeds the AR(1) innovation instead,
\begin{equation*}
\tilde{y}_2[k] = y_2[k] - (1-a_{cov})\,y_2[k-1] = a_{cov}\,s[k],
\end{equation*}
i.e.\ the instantaneous readout $s[k]$ with per-sample variance
$2\,IF_{var}(\bar{a}_w+\xi)^2$ ($+\,d\,Q_{44}$) and (approximately)
independent samples --- the same information rate, in an honest
container.  Verified: $a_{cov}$-invariance scan raw FAIL $3.14\%$
($P[0]$ budget blown $4.645\times$) vs whitened PASS $0.00$--$0.02\%$;
whitening is the production behaviour in both controllers ---
unconditional in \texttt{motion\_control\_law\_5state\_powerlaw.m},
\texttt{y2\_whiten} toggle (default true) in
\texttt{\_5state\_expgain.m} --- the implementation SSOT.  The main text
keeps the raw, white-treated form by decision (user, 2026-07-31) to
mirror the companion documents.

\section*{S9 notes --- Filter equations (not in the main file)}

Both companion documents stop before the filter proper; by decision
(user, 2026-07-31) the filter equations live here.

\begin{align*}
\text{predict:}\quad
&\hat{\mathbf{x}}^-[k+1] = \text{S5 rows}, \qquad
\mathbf{P}^-[k+1] = \mathbf{F}_e[k]\,\mathbf{P}[k]\,\mathbf{F}_e^{\top}[k]
  + \mathbf{Q}[k] \\
\text{update:}\quad
&\mathbf{S}[k] = \mathbf{H}[k]\mathbf{P}^-[k]\mathbf{H}^{\top}[k]
  + \mathbf{R}[k], \qquad
\mathbf{L}[k] = \mathbf{P}^-[k]\mathbf{H}^{\top}[k]\,\mathbf{S}^{-1}[k] \\
&\hat{\mathbf{x}}[k] = \hat{\mathbf{x}}^-[k] + \mathbf{L}[k]\,\mathbf{e}_y[k],
\qquad
\mathbf{P}[k] = (\mathbf{I}-\mathbf{L}\mathbf{H})\,\mathbf{P}^-\,
  (\mathbf{I}-\mathbf{L}\mathbf{H})^{\top}
  + \mathbf{L}\mathbf{R}\mathbf{L}^{\top}
\quad \text{(Joseph)}
\end{align*}
This supplies the correction left open since S6:
$\mathbf{e}[k+1] = \mathbf{F}_e\mathbf{e} - \mathbf{L}\mathbf{H}\mathbf{e}
+ \mathbf{q}$ (the house doc's boxed error state equation).  With diagonal
$\mathbf{R}$ the two channels may be applied as sequential scalar updates
(code practice).  $\mathbf{P}[0]$ from S10; DARE-based initialization per
the implementation SSOT (\texttt{motion\_control\_law\_5state\_expgain\_alg.m}
precedent).

\section*{S10 notes --- Seeds and priors (not in the main file)}

Contract, to be instantiated by the $P[0]$ six-step chain of
\texttt{reference/shared/param\_prior\_rules.md} before any
implementation:
\begin{itemize}
\item Seeds: $\hat{b}_w[0] = 9/8$, $\hat{p}_w[0] = 1$ (the two derived
anchors), $\hat{\bar{w}}_s[0] = 1$ (nominal contact);
$\hat{\bar{a}}_w[0] = \bar{a}(\bar{w}_d[0]-\hat{\bar{w}}_s[0],
\hat{b}_w[0], \hat{p}_w[0])$; $\delta\hat{\bar{w}}_j[0]$ from the first
measurements (or $0$).
\item $P[0]$ widths: $(b_w, p_w)$ funded by the $12.5\%$ anchor tension
via a Form B $\theta_{\mathrm{eff}}$ construction (open task: the Form B
analogue of the expgain $b_{\mathrm{eff}}$ definition);
$\bar{w}_s-1$ width from approach/positioning calibration and the S11
injection protocol; seed errors per-axis via the $A^2$ rule.  $a_o$ is
not a state --- its $2.8\%$ tolerance is exercised by the S11 leakage
test, not carried in $\mathbf{P}$.
\item $Q_\theta = 0$ honesty condition:
$\sup|\theta_{\mathrm{eff}}-\theta_0| \lesssim \sqrt{P[0]}$ (charter
constraint 4).  Collinearity per the four rules: report $\hat{\bar{a}}$
only; never quote $\hat{b}_w$, $\hat{p}_w$ separately; Joseph/UD update;
near-wall separate fit.
\end{itemize}

\section*{S11 notes --- Acceptance criteria (not in the main file)}

Stated before implementation (derivation-workflow rule 3):
\begin{enumerate}
\item Defect-1 fingerprint (S4 risk): adding $y_2$ must \emph{improve}
$\hat{\bar{a}}$ tracking; benchmark the descent peak against the
algebraic-evaluation oracle.
\item $a_{cov}$ invariance: scan $10\times$; whitened must PASS (raw is
expected to FAIL --- that failure is the fingerprint of the raw
container, not of the filter).
\item $P[0]$ budget:
$E[(\hat{x}_\infty-\hat{x}_0)^2] \leq P[0]-P[\infty]$.
\item Shape bound: $\sup|\theta_{\mathrm{eff}}-\theta_0| \lesssim
\sqrt{P[0]}$ with the Form B $\theta_{\mathrm{eff}}$.\\
Tool: adapted \texttt{verify\_shape\_exponent\_bound.m}.
\item $\bar{w}_s$ injection: inject a known wall offset; the filter must
converge to it, on a \emph{moving} trajectory (hold desert, S4/S6 notes).
\item $a_o$ leakage: perturb $a_o$ by $\pm2.8\%$; measure the pull on
$\hat{b}_w$, $\hat{p}_w$, $\hat{\bar{w}}_s$.
\item Truncation bounds: instantiate the S4/S7 $O(\bar{a}'')$ bounds
($\sup|\bar{a}''| = p(p+1)/b^2$, at contact) for the canonical trajectory
class and confirm they are subdominant to the shape bound (item 4).
\end{enumerate}

\end{document}
